\documentclass[10pt,a4paper]{amsart}
\usepackage[margin=19mm]{geometry}
\usepackage{amsmath,amssymb,mathtools}
\usepackage[scr=rsfs,cal=euler]{mathalfa}
\usepackage{xfrac}
\usepackage[unicode,hidelinks,bookmarks=false]{hyperref}
\newcommand{\ie}{i.e.\ }
\newcommand{\cf}{cf.\ }
\newcommand{\resp}{respectively\ }
\newcommand{\Subsection}{\subsection}
\providecommand{\nobreakdash}{\nobreak\hbox{-}}
\setlength{\emergencystretch}{2em}
\multlinegap0pt
\usepackage{amssymb}
\usepackage{mathtools}
\usepackage[shortlabels]{enumitem}
\usepackage{xcolor}
\newtheorem{lemma}{Lemma}
\newtheorem{theorem}{Theorem}
\newtheorem{remark}{Remark}
\newtheorem{proposition}{Proposition}
\newcommand{\LR}{L_{\RR}}
\newcommand{\Mat}[1]{\mathbf{#1}}
\newcommand{\RR}{\mathbb{R}}
\newcommand{\cP}{\mathcal{P}}
\newcommand{\defeq}{\mathrel{\vcentcolon=}}
\newcommand{\vct}[1]{\boldsymbol{#1}}
\title{An algebraic proof of Colombo's difference-power determinant conjecture}
\author{Kun LI\textsuperscript{*}, Li TIE, Peng WANG, Zihan LIU}
\date{Source: arXiv:2608.28274v1 (CC BY 4.0)}
\begin{document}
\maketitle

\noindent\emph{Source and licence.} An algebraic proof of Colombo's difference-power determinant conjecture. Version arXiv:2608.28274v1 (CC BY 4.0), distributed by arXiv under the Creative Commons Attribution 4.0 International licence, http://creativecommons.org/licenses/by/4.0/, as recorded in the arXiv licence metadata for this exact version and shown on the abstract page https://arxiv.org/abs/2608.28274v1. Full licence notice: \texttt{licenses/arxiv-2608.28274-CC-BY-4.0.txt}.\par\smallskip

\noindent\textit{Editorial scope.} Counted unit: the article's proposition prop:odd-nonsingular (source0.tex lines 422--424). The article's complete original proof of it is delivered with it (source0.tex lines 426--482, one \texttt{proof} environment of 57 lines). No mathematics is written, repaired or re-derived: the statement and the proof are the article's own lines, in the article's own notation, and no retained line is substituted or re-flowed.  Retained uncounted as the source-local context the proof needs: lines 340--348; lines 367--371; lines 390--396; lines 408--414.  Nothing was omitted from any retained span, so the package carries no omission record.  No retained line contains a citation, so the wrapper carries no bibliography and no bibliography entry.  No retained line contains a figure environment, an image-inclusion command or drawing code, so no image is delivered and the package declares no asset dependency.  Editorial formatting only, in addition: an AMS \texttt{amsart} preamble that restates the article's own declarations and macros used by the retained lines, namely the environments \texttt{lemma}, \texttt{theorem}, \texttt{remark}, \texttt{proposition} on separate counters, together with the standard packages those lines need.  The retained environments print in this document as Lemma 1, Theorem 1, Remark 1, Proposition 1 in order of appearance, whereas the article prints the same items with the numbering of its own full document.  The complete native source of the article as distributed in the arXiv e-print for this exact version is bundled unchanged as \texttt{sources/208776/source0.tex}.
\begin{lemma}[Pure powers and evaluation]\label{lem:apolar-identities}
For all \(u,v,t\in\RR\) and all \(F\in\cP_d\),
\begin{align}
 \langle L_u^d,L_v^d\rangle_d
 &= (v-u)^d, \label{eq:pure-power-pairing}\\
 \langle L_t^d,F\rangle_d
 &= F(-t,1). \label{eq:evaluation}
\end{align}
\end{lemma}

\begin{lemma}[Classical Vandermonde independence]\label{lem:pure-power-independence}
If \(m\le d+1\) and \(t_1,\ldots,t_m\) are pairwise distinct real
numbers, then \(L_{t_1}^d,\ldots,L_{t_m}^d\) are linearly independent in
\(\cP_d\).
\end{lemma}

\begin{theorem}[Sylvester--Reznick]\label{thm:reznick}
Let \(F\in\RR[X,Y]_d\) be nonzero and not a \(d\)-th power of a real
linear form.  Then
\[
 \tau(F)\le \LR(F).
\]
\end{theorem}

\begin{remark}[Multiplicity matters]\label{rem:multiplicity}
The quantity \(\tau(F)\) counts factors with multiplicity.  Thus an
additional factor increases \(\tau(F)\) even if it coincides with one of
the already prescribed factors.  For example,
\((X-Y)^2(X+Y)\) has only two distinct real projective directions but
has \(\tau=3\), because \(X-Y\) occurs twice.
\end{remark}

\begin{proposition}\label{prop:odd-nonsingular}
The matrix \(\Mat{A}_d(\vct{\lambda})\) is nonsingular.
\end{proposition}

\begin{proof}
Suppose to the contrary that \(\Mat{A}_d(\vct{\lambda})\) is singular,
and choose a nonzero kernel vector
\(\vct{c}=(c_1,\ldots,c_n)^{\mathsf T}\).  Define
\begin{equation}\label{eq:F-kernel}
 F=\sum_{j=1}^n c_jL_{\lambda_j}^d\in\cP_d.
\end{equation}
Since \(n\le d+1\), Lemma~\ref{lem:pure-power-independence} gives
\(F\ne0\).  For every \(r\), Lemma~\ref{lem:apolar-identities} and the
kernel equation give
\[
 F(-\lambda_r,1)
 =\langle L_{\lambda_r}^d,F\rangle_d
 =\sum_{j=1}^n c_j(\lambda_j-\lambda_r)^d
 =-\bigl(\Mat{A}_d\vct{c}\bigr)_r=0.
\]
Consequently,
\[
 X+\lambda_rY\mid F
 \qquad(1\le r\le n).
\]
Indeed, after the linear change \(U=X+\lambda_rY\), \(V=Y\), the value
at \(U=0\) is \(V^dF(-\lambda_r,1)=0\), so every term contains \(U\).

The factors are pairwise nonproportional.  Hence
\begin{equation}\label{eq:Q-divides-F}
 Q(X,Y)\defeq\prod_{r=1}^n(X+\lambda_rY)
 \quad\text{divides}\quad F(X,Y).
\end{equation}
There is therefore a homogeneous \(H\in\RR[X,Y]_{d-n}\) such that
\(F=QH\).  Since \(d>n-1\) and \(d,n\) have opposite parity,
\(d-n\) is a positive odd integer.  Every nonzero real binary form of
odd degree has a real projective zero: on the unit circle,
\(H(-\vct v)=-H(\vct v)\), so the intermediate value theorem along a
semicircle gives a zero.  Thus \(H\) has a real linear factor.
For the concrete example \(H=X^3+Y^3\), the projective zero
\([1:-1]\) corresponds to the factor \(X+Y\).

The factor \(Q\) supplies \(n\) pairwise distinct real linear factors,
and \(H\) supplies one more occurrence, even if it repeats one of them.
Remark~\ref{rem:multiplicity} gives
\begin{equation}\label{eq:tau-lower}
 \tau(F)\ge n+1.
\end{equation}
Because \(Q\) contains at least two nonproportional factors, \(F\) is
not a pure \(d\)-th power.  After zero coefficients are omitted from
\eqref{eq:F-kernel}, the same representation gives
\(\LR(F)\le |\operatorname{supp}(\vct c)|\le n\), where
\[
 \operatorname{supp}(\vct c)=\{j:c_j\ne0\}.
\]
Theorem~\ref{thm:reznick} now yields
\[
 n+1\le\tau(F)\le\LR(F)\le n,
\]
a contradiction.
\end{proof}

\end{document}
